\section{术语}

为了更好地说明 GP-SPI2 的功能，本章使用了以下术语。 

\begin{longtable}[c]{ R{4cm} p{10cm} }

\textbf{主机模式}     & GP-SPI2 用作 SPI 主机，发起 SPI 传输事务。\\ 
\textbf{从机模式}      & GP-SPI2 用作 SPI 从机，当其 CS 被拉低时，与 SPI 主机进行数据传输。\\
\textbf{MISO}       & 主机输入，从机输出。数据从从机发送至主机。\\
\textbf{MOSI}  & 主机输出，从机输入。数据从主机发送至从机。\\ 
\textbf{传输事务} & 一次完整的传输事务：主机拉低从机的 CS 线，开始传输数据，然后再拉高从机的 CS 线。传输事务为原子操作，即不可打断。\\
\textbf{SPI 传输} & SPI 主机与从机完成数据交换的一次完整过程。一次 SPI 传输可以包含一个或多个 SPI 传输事务。\\
\textbf{单次传输} & 在这种传输模式下，仅包含一次传输事务。\\
\textbf{CPU 控制的传输}& 由 CPU 控制，\hyperref[regdesc:SPIW0REG]{SPI\_W0\_REG} \verb+~+ \hyperref[regdesc:SPIW15REG]{SPI\_W15\_REG} 与 SPI 设备之间的数据传输。\\
\textbf{DMA 控制的传输}& 由 DMA 引擎控制，DMA 与 SPI 设备之间的数据传输。\\
\textbf{分段配置传输}& 主机模式下，DMA 控制的数据传输。此类传输包含多个传输事务（分段），每个传输事务均可独立配置。\\
\textbf{从机连续传输} & 从机模式下，DMA 控制的数据传输。此类传输包含多个传输事务（分段）。\\
\textbf{全双工} & 主机与从机之间的发送线和接收线各自独立，发送数据和接收数据同时进行。\\
\textbf{半双工} & 主机和从机只能有一方先发送数据，另一方接收数据。发送数据和接收数据不能同时进行。\\
\textbf{四线全双工} & 四线包括：时钟线、片选线和两条数据线。其中，可使用两条数据线同时发送和接收数据。\\
\textbf{四线半双工} & 四线包括：时钟线、片选线和两条数据线。其中，分时使用两条数据线，不可同时使用。\\
\textbf{三线半双工} & 三线包括：时钟线、片选线和一条数据线。使用数据线分时发送和接收数据。\\
\textbf{1-bit SPI} & 一个时钟周期传输一位数据。\\ 
\textbf{(2-bit) Dual SPI} & 一个时钟周期传输两个数据位。\\
\textbf{Dual Output Read} & Dual SPI 的一种数据模式，一个时钟周期可传输一位命令、或一位地址、或两位数据。\\
\textbf{Dual I/O Read} & Dual SPI 的另外一种数据模式，一个时钟周期可传输一位命令、或两位地址、或两位数据。\\
\textbf{(4-bit) Quad SPI} & 一个时钟周期传输四个数据位。\\
\textbf{Quad Output Read} & Quad SPI 的一种数据模式，一个时钟周期可传输一位命令、或一位地址、或四位数据。\\
\textbf{Quad I/O Read} & Quad SPI 的一种数据模式，一个时钟周期可传输一位命令、或四位地址、或四位数据。\\
\textbf{QPI} & 一个时钟周期可传输四位命令、或四位地址、或四位数据。\\
\end{longtable}
